\documentclass[10pt,a4paper]{amsart}
\usepackage[margin=19mm]{geometry}
\usepackage{amsmath,amssymb,mathtools}
\usepackage[scr=rsfs,cal=euler]{mathalfa}
\usepackage{xfrac}
\usepackage[unicode,hidelinks,bookmarks=false]{hyperref}
\newcommand{\ie}{i.e.\ }
\newcommand{\cf}{cf.\ }
\newcommand{\resp}{respectively\ }
\newcommand{\Subsection}{\subsection}
\providecommand{\nobreakdash}{\nobreak\hbox{-}}
\setlength{\emergencystretch}{2em}
\multlinegap0pt
\usepackage{bm}
\usepackage{bbm}
\usepackage{dsfont}
\usepackage{xspace}
\usepackage{cancel}
\usepackage{mathtools}
\usepackage{amsthm}
\makeatletter
\@ifundefined{Proposition}{\newtheorem{Proposition}{Proposition}}{}
\@ifundefined{Remark}{\newtheorem{Remark}[Proposition]{Remark}}{}
\@ifundefined{Definition}{\newtheorem{Definition}[Proposition]{Definition}}{}
\@ifundefined{Lemma}{\newtheorem{Lemma}[Proposition]{Lemma}}{}
\makeatother
\makeatletter
\newcommand{\Gal}{{\operatorname{Gal}}}
\newcommand{\PP}{{\mathbb P}}
\expandafter\ifx\csname Proof\endcsname\relax
\newenvironment{Proof}{\par\noindent{\sc Proof:}}%
\fi
\expandafter\ifx\csname ProofOf\endcsname\relax
\newenvironment{ProofOf}[1]{\par\noindent{\em Proof of #1.}}%
\fi
\newcommand{\kbar}{{\overline k}}
\makeatother
\title[Visible 2-torsion in the Tate-Shafarevich group of an ellipt]{Visible 2-torsion in the Tate-Shafarevich group of an elliptic curve}
\author{Tom Fisher}
\date{Source: arXiv:2605.31488v1 (CC BY 4.0)}
\begin{document}
\maketitle

\noindent\emph{Source and licence.} Visible 2-torsion in the Tate-Shafarevich group of an elliptic curve. Version arXiv:2605.31488v1 (CC BY 4.0), distributed under the Creative Commons Attribution 4.0 International licence, as recorded in the arXiv licence metadata for this exact version and shown on the abstract page https://arxiv.org/abs/2605.31488v1. Full rights notice: licenses/arxiv-2605.31488-CC-BY-4.0.txt.\par\smallskip

\noindent\textit{Editorial scope.} Counted unit: the Proposition whose source label is \texttt{prop:bs} in the article (source0.tex lines 643--653, source label \texttt{prop:bs}), delivered together with the article's own complete original proof of it (source0.tex lines 655--692, one \texttt{proof} environment of 38 lines). No mathematics is written, repaired or re-derived: statement and proof are the article's own text line for line. Required context retained uncounted: def:Xfg (lines 324--331); rem:8planes (lines 356--361); lem:2planes (lines 380--384). Every definition, display and label of those lines is retained unchanged. Omitted from this package and not redistributed: 1--323, 332--355, 362--379, 385--642, 654--654, 693--1235. Nothing inside a retained excerpt is omitted or altered: every retained body is delivered exactly as the article's lines are written, so no omission receipt of any kind accompanies this package. Citations: the retained excerpts carry no citation command at all, so no bibliography entry of the article is transcribed and no reference is redistributed. Reference adaptation: none was needed and none was made; every reference inside the delivered unit resolves inside it, so no unresolved reference or citation is produced. Printed slips kept literal: no printed word, symbol, hypothesis or number of the counted statement or of its proof was altered. Layout and class: the article's own class is \texttt{amsart} with packages including amssymb, latexsym, comment, url, xy; the wrapper is the lane's pinned \texttt{amsart} editorial wrapper at 10pt on A4, which restates the theorem-like environments the retained text needs and adds the article's own macro definitions that the retained text needs (\texttt{\textbackslash Gal}, \texttt{\textbackslash PP}, \texttt{\textbackslash kbar}), with extra wrapper packages (bm, bbm, dsfont, xspace, cancel, mathtools). The delivered numbering is the unit's own. Assets: no retained span contains a figure environment, a graphics-inclusion command, a PDF-page inclusion, a table or a TikZ picture; no image file is copied into this package, the package declares no asset dependency and no image file is redistributed with it. Licence: version arXiv:2605.31488v1 of this article is distributed under the Creative Commons Attribution 4.0 International licence, as recorded in the arXiv licence metadata for this exact version and shown on the abstract page https://arxiv.org/abs/2605.31488v1. A full rights notice is delivered at licenses/arxiv-2605.31488-CC-BY-4.0.txt. Characters: every retained excerpt is pure ASCII, so the delivered engine, XeLaTeX with its default Latin Modern text font, reports no missing character.
\begin{Definition}
  \label{def:Xfg}
  Let $f$ and $g$ be a pair of non-singular binary quartics with
  the same invariants $I$ and $J$. The {\em associated quadric
  intersection} is 
\[ X_{f,g} = \left\{ \begin{aligned} u_0 u_2 - u_1^2 &= v_0 v_2 - v_1^2 \\
Q_f(u_0,u_1,u_2) &= Q_g(v_0,v_1,v_2) \end{aligned} \right\} \subset \PP^5. \]
\end{Definition}

\begin{Remark}
\label{rem:8planes}
We may change coordinates over $\kbar$ so that $X_{f,g}$ has equations
$\{ x_1 x_2 = x_3 x_4 = x_5 x_6 \} \subset \PP^5$.  It follows that
$X_{f,g}$ contains exactly $8$ planes.
\end{Remark}  

\begin{Lemma}
\label{lem:2planes}
  The planes corresponding to distinct proper
  equivalences intersect only at singular points of $X_{f,g}$.
\end{Lemma}

\begin{Proposition}
  \label{prop:bs}
  Let $X = X_{f,g}$ be as in Definition~\ref{def:Xfg}. Let
  $\Lambda_1, \ldots, \Lambda_4$ be the planes contained in $X$
  corresponding to the proper equivalences between $f$ and $g$, and
  let $H$ be the hyperplane section. Then the linear system
  $|3 H - \Lambda_1 - \Lambda_2 - \Lambda_3 - \Lambda_4|$ defines a
  birational map $\phi : X \, - \!  \to Y$ where $Y \subset \PP^9$ is
  a Brauer-Severi threefold. Moreover, the birational map $\phi$ is
  regular at all smooth points of $X$.
\end{Proposition}

\begin{proof}
  The planes $\Lambda_1, \ldots, \Lambda_4$ are jointly defined over
  $k$, by which we mean that they are permuted by the action of
  $\Gal(\kbar/k)$. Having made this observation, we are free for the
  remainder of the proof to work over $\kbar$.  In particular, we make
  the change of coordinates suggested in Remark~\ref{rem:8planes}.

  We label our coordinates on $\PP^5$ as $x_I = x_{ij}$ where
  $I = \{i,j\}$ runs over the subsets of $\{1,2,3,4\}$ of size 2. By
  Lemma~\ref{lem:2planes} we may assume that
  \[ X = \{ x_{12}x_{34} = x_{13}x_{24} = x_{14}x_{23} \} \subset
    \PP^5 \] and each $\Lambda_i$ is defined by the vanishing of the
  $x_I$ with $i \in I$.  The space of cubic forms vanishing on
  $\Lambda_1 \cup \ldots \cup \Lambda_4$ is spanned by the monomials
  $x_I x_J x_K$ with $\#(I \cup J \cup K) = 4$.  Inside this space, a
  complement to the space of cubic forms vanishing on $X$ is spanned
  by the entries on and above the leading diagonal of the following
  symmetric matrix, where for simplicity we put
  $\alpha = x_{12} x_{34}$.
 \[ \begin{pmatrix}
    x_{12}x_{13}x_{14} & \alpha x_{12} & \alpha x_{13} & \alpha x_{14} \\
    \alpha x_{12} & x_{12}x_{23}x_{24} & \alpha x_{23} & \alpha x_{24} \\
    \alpha x_{13} & \alpha x_{23} & x_{13}x_{23}x_{34} &  \alpha x_{34} \\
    \alpha x_{14} & \alpha x_{24} & \alpha x_{34} & x_{14}x_{24}x_{34}
  \end{pmatrix} \]
  This matrix defines a rational map $\phi : X \, - \!  \to Y$ where
  $Y \subset \PP^9$ is the locus of all $4 \times 4$ symmetric matrices
  of rank $1$.  Let $\psi : \PP^3 \, - \!  \to X$ be given by
 \begin{equation}
 \label{ratparam}
   (x_1: x_2:x_3:x_4) \mapsto (x_{12} : \ldots : x_{34})
   = (x_1x_2 : \ldots : x_3 x_4).
 \end{equation}
 The composite $\phi \circ \psi : \PP^3 \to Y$ is the $2$-uple
 embedding, which is clearly an isomorphism. It follows that $\phi$ is
 birational and $Y$ is a Brauer-Severi variety. The final statement is
 checked by direct calculation.
\end{proof}

\end{document}
